% ============================================================
% content/answers-short.tex —— 2020 全国Ⅰ卷（理综物理部分）· 简明答案速查
% ============================================================

\begin{gknotice}
  选择题答案：14--D，15--B，16--B，17--A，18--C，19--BD，20--AB，21--BC
\end{gknotice}

\section*{一、选择题}

\noindent
\begin{tabular}{@{}l*{8}{c}@{}}
\toprule
题号 & 14 & 15 & 16 & 17 & 18 & 19 & 20 & 21\\
\midrule
答案 & D & B & B & A & C & BD & AB & BC\\
\bottomrule
\end{tabular}

\section*{二、非选择题}

\noindent
\textbf{22.} (1) $O$、$P$；\quad (2) Ⅰ；\quad $50.5~\Omega$；\quad (3) $50.0~\Omega$

\medskip
\noindent
\textbf{23.} (1) 大约相等；\quad (5) $I=m_1gt_{12}$，$\Delta p=m_2\left(\dfrac{d}{\Delta t_2}-\dfrac{d}{\Delta t_1}\right)$；\quad
(6) \SI{0.221}{\newton\second}，\SI{0.212}{\kilogram\metre\per\second}；\quad (7) $4$

\medskip
\noindent
\textbf{24.} (1) \SI{78}{\metre\per\second}；\quad (2) \SI{2.0}{\metre\per\second\squared}，\SI{39}{\second}

\medskip
\noindent
\textbf{25.} (1) $E=\dfrac{mv_0^2}{2qR}$；\quad (2) $v_1=\dfrac{\sqrt{2}}{4}v_0$；\quad (3) $0$ 或 $\dfrac{\sqrt{3}}{2}v_0$

\medskip
\noindent
\textbf{33.} (1) 减小；减小；小于；\quad (2)(i) $\dfrac{2p}{3}$；\quad (ii) $\dfrac{2}{3}$

\medskip
\noindent
\textbf{34.} (1) BCE；\quad (2)(i) $\dfrac{l}{4}$；\quad (ii) $\dfrac{fl}{4}$
